% Archive record 10, line 1.
\hypertarget{part-1-10-0001-conclusion-concept-index}{%
\section{Conclusion \& Concept Index}\label{part-1-10-0001-conclusion-concept-index}}

Part I 的主要目标是建立对 \textbf{Research Problem} 的整体认识，而不是系统地推导 Statistical Learning Theory。前面的数学推导主要用于说明某些现象为什么能够发生，暂时没有必要逐一消化。

\textbf{Conclusion \& Concept Index} 汇集核心定义、数学结论、成立条件及研究意义。具体证明见前文。

\hypertarget{part-1-10-0001-part-i-conclusion-concept-index}{%
\subsection{Part I --- Conclusion \& Concept Index}\label{part-1-10-0001-part-i-conclusion-concept-index}}

\hypertarget{part-1-10-0001-recursive-self-improvement-rsi}{%
\subsection{1. Recursive Self-Improvement (RSI)}\label{part-1-10-0001-recursive-self-improvement-rsi}}

\hypertarget{part-1-10-0001-ux6838ux5fc3ux6982ux5ff5}{%
\subsubsection{1.1 核心概念}\label{part-1-10-0001-ux6838ux5fc3ux6982ux5ff5}}

\begin{readertable}
\begin{longtable}[]{@{}
  >{\raggedright\arraybackslash}p{(\columnwidth - 4\tabcolsep) * \real{0.29}}
  >{\raggedright\arraybackslash}p{(\columnwidth - 4\tabcolsep) * \real{0.55}}
  >{\raggedright\arraybackslash}p{(\columnwidth - 4\tabcolsep) * \real{0.13}}@{}}
\toprule
Concept & 核心含义 & 原文索引 \\
\midrule
\endhead
AI Agent & 由 LLM、Prompt、Memory、Tools、Control Logic 等组成的任务执行系统 & §2.1 \\
Self-Improvement & 系统根据反馈改进自身的方法或结构 & §2.2 \\
Recursive Self-Improvement & 改进后的系统继续参与对自身的改进 & §2.2 \\
Agent Scaffold / Harness & 围绕底层 LLM 构建的程序、工具和控制机制 & §2.1 \\
Improvement Operator & 根据当前版本及历史反馈产生新版本的机制 & §2.2 \\
Trajectory & Agent 在多轮改进中形成的版本演化轨迹 & §3.2 \\
\bottomrule
\end{longtable}
\end{readertable}

RSI 的核心特征可以概括为：

\begin{quote}
The optimizer itself becomes an object of optimization.
\end{quote}

其数学抽象为：
\[A_{k+1}
=
\mathcal I_{A_k}(A_k,F_k,H_k).\]

其中 \(A_k\) 是当前 Agent，\(F_k\) 是实验反馈，\(H_k\) 是历史记录。

\hypertarget{part-1-10-0001-rsi-ux4e0eux5176ux4ed6ux4f18ux5316ux65b9ux6cd5ux7684ux5173ux7cfb}{%
\subsubsection{1.2 RSI 与其他优化方法的关系}\label{part-1-10-0001-rsi-ux4e0eux5176ux4ed6ux4f18ux5316ux65b9ux6cd5ux7684ux5173ux7cfb}}

\begin{readertable}
\begin{longtable}[]{@{}
  >{\raggedright\arraybackslash}p{(\columnwidth - 6\tabcolsep) * \real{0.16}}
  >{\raggedright\arraybackslash}p{(\columnwidth - 6\tabcolsep) * \real{0.19}}
  >{\raggedright\arraybackslash}p{(\columnwidth - 6\tabcolsep) * \real{0.31}}
  >{\raggedright\arraybackslash}p{(\columnwidth - 6\tabcolsep) * \real{0.29}}@{}}
\toprule
& ML Training & Hyperparameter Tuning & RSI \\
\midrule
\endhead
优化对象 & Parameters & Configurations & Agent / Method \\
搜索空间 & 参数空间 & 配置空间 & 程序与方法空间 \\
优化机制 & 通常固定 & 通常固定 & 可能被自身修改 \\
自适应性 & 存在 & 存在 & 存在，且可能更开放 \\
自指性 & 通常没有 & 通常没有 & 核心特征 \\
\bottomrule
\end{longtable}
\end{readertable}

\textbf{重要观点：} RSI 并没有创造一种全新的统计学问题，而是可能通过开放式搜索、自适应反馈及自我修改，放大已有的统计学习问题。

证明与解释索引：§2--4。

\hypertarget{part-1-10-0001-ux6838ux5fc3ux6570ux5b66ux7ed3ux8bba}{%
\subsection{2. 核心数学结论}\label{part-1-10-0001-ux6838ux5fc3ux6570ux5b66ux7ed3ux8bba}}

\hypertarget{part-1-10-0001-conclusion-1-in-sample-improvement-does-not-imply-out-of-sample-improvement}{%
\subsubsection{Conclusion 1 --- In-Sample Improvement does not imply Out-of-Sample Improvement}\label{part-1-10-0001-conclusion-1-in-sample-improvement-does-not-imply-out-of-sample-improvement}}

在固定评估环境和允许保留原版本的 Greedy Selection 机制下：
\[\boxed{
\widehat J_{\mathrm{IS}}(A_{k+1})
\geq
\widehat J_{\mathrm{IS}}(A_k)
}\]

但是：
\[\boxed{
\Delta J_{\mathrm{IS}}>0
\quad\not\Rightarrow\quad
\Delta J_{\mathrm{OOS}}>0
}\]

\textbf{数学性质：} 前者是特定选择规则下的严格结论，后者是不成立的逻辑蕴含。

\textbf{研究意义：} IS Score 的单调增长可能仅仅是优化规则的直接结果，而不是 Agent 真实性能持续提升的证据。

Proof Index：§3.1 --- Generate--Evaluate--Select Loop

\hypertarget{part-1-10-0001-conclusion-2-unbiased-evaluation-does-not-imply-unbiased-model-selection}{%
\subsubsection{Conclusion 2 --- Unbiased Evaluation does not imply Unbiased Model Selection}\label{part-1-10-0001-conclusion-2-unbiased-evaluation-does-not-imply-unbiased-model-selection}}

对于事先固定且独立于评估数据的模型 \(A\)，在评估样本同分布且期望存在的条件下：
\[\mathbb E_D[\widehat J_D(A)]
=
J(A).\]

但是，当模型本身依赖于评估数据时：
\[\boxed{
\mathbb E_D
\left[
\widehat J_D(A(D))-J(A(D))
\right]
\text{ 不一定为 }0
}\]

\textbf{数学性质：} 严格的统计学观点。

\textbf{研究意义：} 即使 Evaluator 对每一个事先固定的 Agent 都是无偏的，经过 Evaluator 反复筛选得到的 Agent，其观察得分仍然可能存在系统性偏差。

这里的关键是 \textbf{Data-dependent Model Selection}。

Proof Index：§5.1 --- Population Performance and Empirical Performance

\hypertarget{part-1-10-0001-conclusion-3-model-selection-introduces-optimistic-bias}{%
\subsubsection{Conclusion 3 --- Model Selection introduces Optimistic Bias}\label{part-1-10-0001-conclusion-3-model-selection-introduces-optimistic-bias}}

假设所有候选模型的真实性能相同：
\[J(A_1)=\cdots=J(A_m)=\mu.\]

观察得分为：
\[\widehat J_D(A_j)=\mu+\varepsilon_j,
\qquad
\mathbb E[\varepsilon_j]=0.\]

选择得分最高的模型后：
\[\boxed{
\mathbb E[\widehat J_D(A_{j^*})]
\geq\mu
}\]

但其真实性能仍然为：
\[J(A_{j^*})=\mu.\]

\textbf{数学性质：} 严格结论，利用 Maximum 的凸性与 Jensen's Inequality 即可证明，不要求各评估误差相互独立。

\textbf{研究意义：} 搜索并选择最优模型，本身就可能产生虚假的性能提升。

相关概念：Selection Bias、Winner's Curse、Multiple Comparisons。

Proof Index：§5.2 --- Selection Bias

\hypertarget{part-1-10-0001-conclusion-4-selection-bias-increases-with-search-budget-and-evaluation-noise}{%
\subsubsection{Conclusion 4 --- Selection Bias increases with Search Budget and Evaluation Noise}\label{part-1-10-0001-conclusion-4-selection-bias-increases-with-search-budget-and-evaluation-noise}}

在独立 Gaussian Evaluation Error 的假设下：
\[\varepsilon_j
\overset{\mathrm{i.i.d.}}{\sim}
\mathcal N(0,\sigma^2),\]

当候选模型数量 \(m\) 足够大时：
\[\boxed{
\mathbb E[\max_j\varepsilon_j]
\sim
\sigma\sqrt{2\log m}
}\]

进一步，在独立样本均值等适当条件下：
\[\sigma\propto\frac1{\sqrt n}.\]

因此 Selection Bias 具有如下尺度关系：
\[\boxed{
\text{Selection Bias}
=
O\left(\sqrt{\frac{\log m}{n}}\right)
}\]

\bookgraphic{0.615}{part-1-10-0001-004-01}{}

n = 100

n = 500

n = 2,000

基于 Gaussian Maximum 渐近近似的理论示意，假设单样本误差方差为 1。

\textbf{数学性质：} 有条件的渐近结论，不是适用于所有 RSI 系统的普遍定理。

\textbf{研究意义：} Agent 搜索能力的增强不一定总是有益的。搜索预算越大、评估数据越少、评估噪声越大，模型选择的可靠性就越值得关注。

Proof Index：§5.2 --- Gaussian Maximum and Search Budget

\hypertarget{part-1-10-0001-conclusion-5-score-inflation-ux4e0e-performance-decay-ux662fux4e0dux540cux73b0ux8c61}{%
\subsubsection{Conclusion 5 --- Score Inflation 与 Performance Decay 是不同现象}\label{part-1-10-0001-conclusion-5-score-inflation-ux4e0e-performance-decay-ux662fux4e0dux540cux73b0ux8c61}}

\textbf{Score Inflation}

\bookgraphic{0.445}{part-1-10-0001-006-01}{}

\[\Delta J_{\mathrm{IS}}>0\]

\[\Delta J_{\mathrm{OOS}}=0\]

观察表现增长，真实性能不变。

\textbf{Performance Decay}

\bookgraphic{0.445}{part-1-10-0001-006-02}{}

\[\Delta J_{\mathrm{IS}}>0\]

\[\Delta J_{\mathrm{OOS}}<0\]

观察表现增长，样本外性能下降。

IS Performance

Population / OOS Performance

\textbf{重要逻辑关系：}

\[\text{Selection Bias}
\not\Rightarrow
\text{Performance Decay}.\]

Selection Bias 可以解释观察性能的高估，却不能单独证明真实性能下降。

而 X Program \#161 真正希望研究的是第二种现象及其产生机制。

Explanation Index：§5.3

\hypertarget{part-1-10-0001-conclusion-6-adding-noise-features-can-strictly-damage-generalization}{%
\subsubsection{Conclusion 6 --- Adding Noise Features can strictly damage Generalization}\label{part-1-10-0001-conclusion-6-adding-noise-features-can-strictly-damage-generalization}}

通过前面构建的 Linear Regression Example：
\[Y=\beta X_1+\varepsilon,\]

其中 \(X_2\) 是独立的无效特征。 设：
\[f_0(X)=\beta X_1,\]

\[f_1(X)=\beta X_1+\hat bX_2.\]

在此前设定的模型假设下，我们严格证明了：
\[\boxed{
\begin{aligned}
\widehat R_D(f_1)
&\leq\widehat R_D(f_0),\\
R(f_1)&\geq R(f_0).
\end{aligned}
}\]

尤其是：
\[R(f_1)-R(f_0)=\hat b^2.\]

\textbf{数学性质：} 严格结论。在此前连续分布及非退化假设下，两个不等式几乎必然严格成立。

\textbf{研究意义：} 即使不存在金融市场变化、复杂模型或 LLM，单纯增加无效特征也足以造成真实的泛化性能下降。

这个例子展示的是经典 Empirical Risk Overfitting，而非单独证明 Adaptive Overfitting。

Proof Index：§5.4 --- A Rigorous Regression Example

\hypertarget{part-1-10-0001-conclusion-7-repeated-adaptive-evaluation-compromises-holdout-independence}{%
\subsubsection{Conclusion 7 --- Repeated Adaptive Evaluation compromises Holdout Independence}\label{part-1-10-0001-conclusion-7-repeated-adaptive-evaluation-compromises-holdout-independence}}

RSI 中存在如下信息传递关系：
\[D
\longrightarrow H_k
\longrightarrow A_{k+1}.\]

其中 \(D\) 为反复使用的评估数据集。

\textbf{Evaluation Dataset \(D\)}

↓

\textbf{Evaluation Feedback}

Score、Error Analysis、Failure Cases

↓

\textbf{Experimental History \(H_k\)}

↓

\textbf{Next Agent \(A_{k+1}\)}

间接依赖于 Evaluation Dataset

因此，Agent 即使没有直接读取评估数据的标签，也可能通过反复获取反馈逐渐适应该数据集。

\textbf{研究意义：} Validation Set 一旦反复参与模型改进，就不能继续被无条件地视为独立评估集。

这属于 Adaptive Data Analysis 的核心问题。

需要注意，自适应依赖会使经典泛化保证不再自动成立，但并不意味着所有自适应分析都必然过拟合。

Explanation Index：§6 --- Adaptive Overfitting

\hypertarget{part-1-10-0001-conclusion-8-statistical-overfitting-ux4e0e-distribution-shift-ux5fc5ux987bux533aux5206}{%
\subsubsection{Conclusion 8 --- Statistical Overfitting 与 Distribution Shift 必须区分}\label{part-1-10-0001-conclusion-8-statistical-overfitting-ux4e0e-distribution-shift-ux5fc5ux987bux533aux5206}}

设 \(P\) 为开发阶段的数据分布，\(Q\) 为目标环境的数据分布。

则：
\[\boxed{
\begin{aligned}
\widehat J_D(A)-J_Q(A)
={}&
\underbrace{
\widehat J_D(A)-J_P(A)
}_{\text{Statistical Generalization Gap}}\\
&+
\underbrace{
J_P(A)-J_Q(A)
}_{\text{Distribution Shift Gap}}
\end{aligned}
}\]

\textbf{数学性质：} 严格的代数恒等式。

\textbf{研究意义：} 观察到金融模型的 OOS Performance 下降，并不足以说明下降来自 RSI Overfitting。

我们需要区分模型是否过度适应有限样本，以及市场的数据生成机制是否发生变化。

两个差值均可能为正或负，且该分解本身并不能保证我们能够从实验数据中分别识别它们。

Proof Index：§7.2 --- Non-stationarity and Distribution Shift

\hypertarget{part-1-10-0001-ux4e3aux4ec0ux4e48-finance-ux5c24ux5176ux5bb9ux6613ux51faux73b0ux8fd9ux4e9bux95eeux9898}{%
\subsection{3. 为什么 Finance 尤其容易出现这些问题？}\label{part-1-10-0001-ux4e3aux4ec0ux4e48-finance-ux5c24ux5176ux5bb9ux6613ux51faux73b0ux8fd9ux4e9bux95eeux9898}}

这一部分目前只需要掌握各个因素与研究问题的逻辑联系。

\begin{readertable}
\begin{longtable}[]{@{}
  >{\raggedright\arraybackslash}p{(\columnwidth - 2\tabcolsep) * \real{0.41}}
  >{\raggedright\arraybackslash}p{(\columnwidth - 2\tabcolsep) * \real{0.55}}@{}}
\toprule
Financial Property & Consequence \\
\midrule
\endhead
Low Signal-to-Noise Ratio & 难以区分真实预测信号和偶然相关性 \\
Non-stationarity & 历史规律可能在未来失效 \\
Serial Dependence & 常规独立样本统计推断可能不再有效 \\
Volatility Clustering & 不同时间段的评估不确定性不同 \\
Heavy-tailed Returns & 极端值可能显著影响评价指标 \\
Backtest Overfitting & 大规模策略搜索可能产生虚假历史优势 \\
Look-Ahead Bias & 使用决策时尚不可获得的信息 \\
Transaction Costs & 微小的历史性能提升可能没有实际价值 \\
Benchmark Contamination & LLM 可能已经接触过部分历史评估信息 \\
\bottomrule
\end{longtable}
\end{readertable}

其中最重要的三个因素是：
\[\boxed{
\begin{aligned}
&\text{Low SNR}\\
&\text{Adaptive Model Selection}\\
&\text{Distribution Shift}
\end{aligned}
}\]

前两个因素可能共同加剧过拟合；第三个因素则可能独立造成性能退化，并使其原因更加难以识别。

Explanation Index：§7 --- Financial AI

\hypertarget{part-1-10-0001-x-program-161-research-methodology}{%
\subsection{4. X Program \#161 --- Research Methodology}\label{part-1-10-0001-x-program-161-research-methodology}}

这部分是 Part I 中与实际科研任务联系最紧密的内容。

\hypertarget{part-1-10-0001-trajectory-analysis}{%
\subsubsection{4.1 Trajectory Analysis}\label{part-1-10-0001-trajectory-analysis}}

对于每一轮 Agent，记录：
\[\mathcal T=
\left\{
(A_k,S_k^{\mathrm{IS}},S_k^{\mathrm{OOS}},r_k)
\right\}_{k=0}^{K}.\]

计算经验 Generalization Gap：
\[\boxed{
G_k=S_k^{\mathrm{IS}}-S_k^{\mathrm{OOS}}
}\]

同时需要区分 Chronological Trajectory、Incumbent Trajectory 与 Lineage。

Index：§3.2、§8.1

\hypertarget{part-1-10-0001-decay-detection}{%
\subsubsection{4.2 Decay Detection}\label{part-1-10-0001-decay-detection}}

核心任务：识别 OOS Performance 从改善转向退化的位置。

\bookgraphic{0.786}{part-1-10-0001-015-01}{}

假想的 RSI Performance Trajectory。

候选数学工具：Segmented Regression。

\[S_k^{\mathrm{OOS}}
=
\alpha+\beta k+\gamma(k-\tau)_+
+\eta_k.\]

若：
\[\beta>0,
\qquad
\beta+\gamma<0,\]

则该模型描述了先改善、后退化的趋势。

需要区分 Observed Peak、Change Point 与 Sustained Decay，而不能直接将观察最大值视为真实变点。

Index：§8.2

\hypertarget{part-1-10-0001-revision-classification}{%
\subsubsection{4.3 Revision Classification}\label{part-1-10-0001-revision-classification}}

对每一次修改计算：
\[\begin{aligned}
\Delta_k^{\mathrm{IS}}
&=S_k^{\mathrm{IS}}-S_{k-1}^{\mathrm{IS}},\\
\Delta_k^{\mathrm{OOS}}
&=S_k^{\mathrm{OOS}}-S_{k-1}^{\mathrm{OOS}}.
\end{aligned}\]

以及：
\[\boxed{
\Delta G_k
=
\Delta_k^{\mathrm{IS}}
-
\Delta_k^{\mathrm{OOS}}
}\]

\bookgraphic{0.444}{part-1-10-0001-018-01}{重点分析提高 IS 表现却降低 OOS 表现的 Revision。}

候选 Revision Categories：

\begin{readertable}
\begin{longtable}[]{@{}
  >{\raggedright\arraybackslash}p{(\columnwidth - 2\tabcolsep) * \real{0.51}}
  >{\raggedright\arraybackslash}p{(\columnwidth - 2\tabcolsep) * \real{0.45}}@{}}
\toprule
Category & 修改对象 \\
\midrule
\endhead
Feature Engineering & 特征构造 \\
Model Complexity & 模型结构与复杂度 \\
Search Expansion & 搜索范围及预算 \\
Memory Modification & 历史反馈的利用方式 \\
Evaluation Modification & 评分和选择机制 \\
Data Processing & 数据处理流程 \\
Robustness Mechanism & 正则化与稳健性机制 \\
Infrastructure & 执行效率及工具接口 \\
\bottomrule
\end{longtable}
\end{readertable}

\textbf{研究原则：} 分类统计只能发现关联，进一步的因果解释需要 Rollback Experiment、Ablation Study 等受控实验。

Index：§8.3

\hypertarget{part-1-10-0001-overfitting-mitigation}{%
\subsubsection{4.4 Overfitting Mitigation}\label{part-1-10-0001-overfitting-mitigation}}

\begin{readertable}
\begin{longtable}[]{@{}
  >{\raggedright\arraybackslash}p{(\columnwidth - 2\tabcolsep) * \real{0.46}}
  >{\raggedright\arraybackslash}p{(\columnwidth - 2\tabcolsep) * \real{0.50}}@{}}
\toprule
Mechanism & 基本思想 \\
\midrule
\endhead
Regularization & 限制不必要的复杂度 \\
Early Stopping & 在后期退化前停止优化 \\
Independent Validation & 减少对固定评估集的适应 \\
Rolling Validation & 检验不同时间窗口下的稳健性 \\
Search Budget Control & 减少大规模搜索的选择偏差 \\
Statistical Acceptance Test & 不接受仅由随机波动产生的提升 \\
Revision Rollback & 撤销有害的修改 \\
Robust Objective & 将跨环境稳定性纳入优化目标 \\
\bottomrule
\end{longtable}
\end{readertable}

例如，将传统接受条件：
\[\Delta_k^{\mathrm{IS}}>0\]

修改为：
\[\Delta_k^{\mathrm{IS}}
>
c\cdot
\widehat{\operatorname{SE}}
(\Delta_k^{\mathrm{IS}}).\]

不过，反复执行统计检验本身也可能引入 Multiple Testing 和 Adaptive Overfitting，因此这并不是具有无条件统计保证的解决方案。

Index：§8.4

以上 Detection、Taxonomy 和 Mitigation 方法均为我们目前提出的候选研究方案，不代表课题组已经采用这些方法。

\hypertarget{part-1-10-0001-literature-index}{%
\subsection{5. Literature Index}\label{part-1-10-0001-literature-index}}

暂时不展开论文内容，仅保留文献与研究问题之间的映射关系。

\begin{readertable}
\begin{longtable}[]{@{}
  >{\raggedright\arraybackslash}p{(\columnwidth - 4\tabcolsep) * \real{0.30}}
  >{\raggedleft\arraybackslash}p{(\columnwidth - 4\tabcolsep) * \real{0.40}}
  >{\raggedright\arraybackslash}p{(\columnwidth - 4\tabcolsep) * \real{0.30}}@{}}
\toprule
Paper & Year & Research Connection \\
\midrule
\endhead
\href{https://www.jmlr.org/papers/v11/cawley10a.html}{On Over-fitting in Model Selection and Subsequent Selection Bias in Performance Evaluation} & 2010 & Model Selection Bias \\
\href{https://doi.org/10.1126/science.aaa9375}{The Reusable Holdout} & 2015 & Adaptive Data Analysis \\
The Probability of Backtest Overfitting & 2017 & Financial Evaluation \\
\href{https://arxiv.org/abs/2310.02304}{Self-Taught Optimizer (STOP)} & 2024 & RSI 的数学与程序结构 \\
\href{https://arxiv.org/abs/2505.22954}{Darwin Gödel Machine} & 2025/26 & Agent Evolution、Archive \\
\href{https://arxiv.org/abs/2505.15155}{R\&D-Agent-Quant} & 2025 & Automated Financial Research \\
\href{https://arxiv.org/abs/2609.26457}{Recursive self-improvement of AI research agents} & 2026 & RSI 与跨任务泛化 \\
\bottomrule
\end{longtable}
\end{readertable}

完整作者信息、出版状态和论文介绍索引至 §9。这里沿用前面的文献清单，未进行新一轮文献核验。

\hypertarget{part-1-10-0001-part-i-knowledge-checklist}{%
\subsection{6. Part I --- Knowledge Checklist}\label{part-1-10-0001-part-i-knowledge-checklist}}

我们暂时不要求能够独立复现所有数学推导，只要求能够理解并解释以下结论。

\textbf{Part I 掌握程度}

\begin{itemize}
\tightlist
\item
  RSI 的定义及其与传统优化方法的区别
\end{itemize}

§2--4

\begin{itemize}
\tightlist
\item
  IS Performance 增长不保证 OOS Performance 增长
\end{itemize}

§3.1

\begin{itemize}
\tightlist
\item
  固定模型的无偏评估不保证模型选择的无偏性
\end{itemize}

§5.1

\begin{itemize}
\tightlist
\item
  Selection Bias 及其与 Search Budget 的关系
\end{itemize}

§5.2

\begin{itemize}
\tightlist
\item
  Score Inflation 与 Performance Decay 的区别
\end{itemize}

§5.3

\begin{itemize}
\tightlist
\item
  无效特征如何造成真实性能退化
\end{itemize}

§5.4

\begin{itemize}
\tightlist
\item
  Adaptive Feedback 如何破坏 Holdout Independence
\end{itemize}

§6

\begin{itemize}
\tightlist
\item
  Statistical Overfitting 与 Distribution Shift 的区别
\end{itemize}

§7

\begin{itemize}
\tightlist
\item
  Trajectory Analysis 与 Decay Detection
\end{itemize}

§8.1--2

\begin{itemize}
\tightlist
\item
  Revision Classification 与 Overfitting Mitigation
\end{itemize}

§8.3--4

\hypertarget{part-1-10-0001-ux6700ux7ec8ux7684-conceptual-framework}{%
\subsection{7. 最终的 Conceptual Framework}\label{part-1-10-0001-ux6700ux7ec8ux7684-conceptual-framework}}

如果我们将整个 Part I 压缩到只剩一条主线，那就是：

\textbf{Recursive Self-Improvement}

系统不断改进自身方法

↓ \textbf{Repeated \& Adaptive Model Selection}

产生大量候选版本并反复获取评估反馈

↓ \textbf{Selection Bias \& Adaptive Overfitting}

可能将随机规律误认为有效信号

↓ \textbf{Generalization Gap Expansion}

开发集与样本外表现逐渐分离

↓ \textbf{Potential Out-of-Sample Performance Decay}

样本外性能可能真正下降

\textbf{Detection}

何时退化

\textbf{Diagnosis}

为何退化

\textbf{Mitigation}

如何缓解

其中，金融环境中的 Low SNR、Non-stationarity 和时间依赖性进一步增加了研究难度。

后面进入 Part II 时，我建议调整一下教学颗粒度：\textbf{每次只推进一个相对独立的理论模块，采用 Definition → Key Results → Mathematical Principles → Research Implications 的结构。} 每个模块结束后单独整理 Conclusion Index，避免再次出现大量概念与证明同时堆积的情况。

Part I 到这里就可以暂时结束，不必为了完全掌握每一个证明而阻塞后续学习。
